% Generated draft: requires mathematical and source review.
\documentclass[11pt]{article}
\usepackage[margin=25mm]{geometry}
\usepackage{fontspec}
\setmainfont{Latin Modern Roman}
\usepackage{amsmath,amssymb,mathtools,unicode-math}
\setmathfont{Latin Modern Math}
\usepackage{enumitem,needspace,xurl,hyperref}
\setlength{\parindent}{0pt}
\setlength{\parskip}{5pt}
\setlength{\emergencystretch}{3em}
\newcommand{\N}{\mathbb{N}}
\newcommand{\Z}{\mathbb{Z}}
\newcommand{\Q}{\mathbb{Q}}
\newcommand{\R}{\mathbb{R}}
\newcommand{\C}{\mathbb{C}}
\newcommand{\F}{\mathbb{F}}
\newcommand{\A}{\mathbb{A}}
\newcommand{\PP}{\mathbb{P}}
\newcommand{\E}{\mathbb{E}}
\newcommand{\eps}{\varepsilon}
\newcommand{\dd}{\,\mathrm{d}}
\begin{document}
\section*{UnsolvedMath ID: 30004697}
\textbf{Source problem code:} catalogue entry
\subsection{Definitions and mathematical statement}
The supplied record concerns a historical inquiry about mathematical texts, rather than a conjecture with specified mathematical objects, hypotheses, and conclusion. The titles \emph{The Nine Chapters} and Yang Hui's \emph{Reclassification} refer to the works named in the source; their exact editions, relevant passages, and problem lists are not supplied.

\textbf{Source inquiry.} Investigate the following questions, retaining the historical premises as supplied claims requiring documentary verification:
\begin{enumerate}
\item Why did Yang Hui and Jia Xian add nine problems to \emph{The Nine Chapters}?
\item Why did Yang Hui choose the added problems rather than the original problems in his \emph{Reclassification}?
\item Why did Yang Hui change their order?
\end{enumerate}
The record does not identify the nine problems, specify the original and revised orders, explain the respective roles of Yang Hui and Jia Xian, or unambiguously identify the collection referred to by ``their'' in the third question. Consequently, a more precise historical formulation requires additional textual evidence. No mathematical proposition or criterion for a definitive historical explanation is specified, so a quantified mathematical statement cannot be extracted without introducing unsupported content.
\subsection{Short English statement}
Explain the historical reasons for the nine added problems attributed to Yang Hui and Jia Xian, Yang Hui's selection of added problems for his \emph{Reclassification}, and his reordering of problems. The relevant texts, problem identities, and ordering comparison require clarification. This is a historical research question requiring human review.
\subsection{Sources}
\begin{itemize}
\item Local registry: \texttt{data/unsolvedmath/problems.json}, permanent catalogue entry, source problem code \texttt{catalogue entry}, entitled ``Yang Hui's Reclassification of Problems from The Nine Chapters.'' The supplied record belongs to the Oberwolfach Reports Open Problems collection, described as extracting open and partially solved questions from workshop discussions. No individual proposer is supplied; the recorded proposed year is 2021.
\item Supplied report citation: \emph{Mathematics and its Ancient Classics Worldwide: Translations, Appropriations, Reconstructions, Roles} (2022), \emph{Oberwolfach Reports}, pp.\ 1347--1406, DOI \texttt{10.4171/owr/2021/26}. Supplied source URL and literature reference labelled ``Oberwolfach report'': \url{https://doi.org/10.4171/owr/2021/26}. This report citation is supplied provenance; its contents have not been independently inspected here.
\item Literature assessment dated exactly \texttt{2026-08-22}: the supplied curation labels the question ``open'' and states that the motivations for Yang Hui's additions, selections, and reordering lack a definitive explanation in the cited primary scholarship. This is supplied dated curation, not an independently verified literature assessment. The background explicitly calls for independent verification of time-sensitive status and cited claims; no separate primary textual editions or scholarship citations are identified in the record.
\end{itemize}
\end{document}
